% GENERATED FILE. Edit canonical lessons, then run npm run generate:books.
% canonical-source-hash: fnv1a64:4b6d722313760745
% canonical-lessons: JA-C94-konomu, JA-C94-konomi, JA-C94-shumi, JA-C94-tanoshimi, JA-C94-ai

\chapter{To like, Liking, Hobby, Something to look forward to, Love}
\label{ch:ja-to-like-liking-hobby-something-to-look-forward-to-love}
\hlchaptermodality{94}

\begin{chapteropening}
\textbf{By the end of this chapter:} \emph{I can say to like, a liking, a hobby, something to look forward to and love.}

Complete the last lesson of chapter 94: I can say to like, a liking, a hobby, something to look forward to and love.
\end{chapteropening}

\section[\texorpdfstring{\ja{このむ} (konomu)}{このむ (konomu)}]{\ja{このむ} (konomu) — to like, to prefer}

\label{lesson:JA-C94-konomu}



\begin{quote}
Before the new one: say the Japanese for a candidate, then the Japanese for a plan.
\end{quote}

\subsection*{You'll want to know: \ja{このむ}}
\textbf{\ja{このむ}} — \emph{konomu} — \textquotedblleft{}to like, to prefer\textquotedblright{}.

\begin{grammarlens}[title={What you've built}]
One more word for this chapter.
\end{grammarlens}

\subsection*{Your turn}
\begin{itemize}
\raggedright
  \item \emph{Say it:} \emph{konomu}
  \item \emph{Say it:} \emph{konomu}, once more
  \item \emph{Say it:} say \emph{konomu}
  \item \emph{Recall:} say the Japanese for a candidate, then the Japanese for a plan, then say \emph{konomu} again
\end{itemize}

\begin{tcolorbox}[breakable,colback=teal!4,colframe=teal!35!black,title={Before you move on}]
What does \ja{このむ} mean? (\textquotedblleft{}To like, to prefer\textquotedblright{}.) Say it once more.
\end{tcolorbox}
\section[\texorpdfstring{\ja{このみ} (konomi)}{このみ (konomi)}]{\ja{このみ} (konomi) — a liking, taste}

\label{lesson:JA-C94-konomi}



\begin{quote}
Before the new one: say the Japanese for a plan, then the Japanese for to like.
\end{quote}

\subsection*{You'll want to know: \ja{このみ}}
\textbf{\ja{このみ}} — \emph{konomi} — \textquotedblleft{}a liking, taste\textquotedblright{}.

\begin{grammarlens}[title={What you've built}]
One more word for this chapter.
\end{grammarlens}

\subsection*{Your turn}
\begin{itemize}
\raggedright
  \item \emph{Say it:} \emph{konomi}
  \item \emph{Say it:} \emph{konomi}, once more
  \item \emph{Say it:} say \emph{konomi}
  \item \emph{Recall:} say the Japanese for a plan, then the Japanese for to like, then say \emph{konomi} again
\end{itemize}

\begin{tcolorbox}[breakable,colback=teal!4,colframe=teal!35!black,title={Before you move on}]
What does \ja{このみ} mean? (\textquotedblleft{}A liking, taste\textquotedblright{}.) Say it once more.
\end{tcolorbox}
\section[\texorpdfstring{\ja{しゅみ} (shumi)}{しゅみ (shumi)}]{\ja{しゅみ} (shumi) — a hobby}

\label{lesson:JA-C94-shumi}



\begin{quote}
Before the new one: say the Japanese for to like, then the Japanese for a liking.
\end{quote}

\subsection*{You'll want to know: \ja{しゅみ}}
\textbf{\ja{しゅみ}} — \emph{shumi} — \textquotedblleft{}a hobby\textquotedblright{}.

\begin{grammarlens}[title={What you've built}]
One more word for this chapter.
\end{grammarlens}

\subsection*{Your turn}
\begin{itemize}
\raggedright
  \item \emph{Say it:} \emph{shumi}
  \item \emph{Say it:} \emph{shumi}, once more
  \item \emph{Say it:} say \emph{shumi}
  \item \emph{Recall:} say the Japanese for to like, then the Japanese for a liking, then say \emph{shumi} again
\end{itemize}

\begin{tcolorbox}[breakable,colback=teal!4,colframe=teal!35!black,title={Before you move on}]
What does \ja{しゅみ} mean? (\textquotedblleft{}A hobby\textquotedblright{}.) Say it once more.
\end{tcolorbox}
\section[\texorpdfstring{\ja{たのしみ} (tanoshimi)}{たのしみ (tanoshimi)}]{\ja{たのしみ} (tanoshimi) — something to look forward to}

\label{lesson:JA-C94-tanoshimi}



\begin{quote}
Before the new one: say the Japanese for a liking, then the Japanese for a hobby.
\end{quote}

\subsection*{You'll want to know: \ja{たのしみ}}
\textbf{\ja{たのしみ}} — \emph{tanoshimi} — \textquotedblleft{}something to look forward to\textquotedblright{}.

\begin{grammarlens}[title={What you've built}]
One more word for this chapter.
\end{grammarlens}

\subsection*{Your turn}
\begin{itemize}
\raggedright
  \item \emph{Say it:} \emph{tanoshimi}
  \item \emph{Say it:} \emph{tanoshimi}, once more
  \item \emph{Say it:} say \emph{tanoshimi}
  \item \emph{Recall:} say the Japanese for a liking, then the Japanese for a hobby, then say \emph{tanoshimi} again
\end{itemize}

\begin{tcolorbox}[breakable,colback=teal!4,colframe=teal!35!black,title={Before you move on}]
What does \ja{たのしみ} mean? (\textquotedblleft{}Something to look forward to\textquotedblright{}.) Say it once more.
\end{tcolorbox}
\section[\texorpdfstring{\ja{あい} (ai)}{あい (ai)}]{\ja{あい} (ai) — love}

\label{lesson:JA-C94-ai}



\begin{quote}
Before the new one: say the Japanese for a hobby, then the Japanese for something to look forward to.
\end{quote}

\subsection*{You'll want to know: \ja{あい}}
\textbf{\ja{あい}} — \emph{ai} — \textquotedblleft{}love\textquotedblright{}.

\begin{grammarlens}[title={What you've built}]
One more word for this chapter.
\end{grammarlens}

\subsection*{Your turn}
\begin{itemize}
\raggedright
  \item \emph{Say it:} \emph{ai}
  \item \emph{Say it:} \emph{ai}, once more
  \item \emph{Say it:} say \emph{ai}
  \item \emph{Recall:} say the Japanese for a hobby, then the Japanese for something to look forward to, then say \emph{ai} again
\end{itemize}

\begin{tcolorbox}[breakable,colback=teal!4,colframe=teal!35!black,title={Before you move on}]
What does \ja{あい} mean? (\textquotedblleft{}Love\textquotedblright{}.) Say it once more.
\end{tcolorbox}
